\section{Literature Survey and Technology Stack}

In recent years, web application development and artificial intelligence have progressed rapidly. Standard web portals have evolved from static HTML/CSS pages into highly dynamic Single Page Applications (SPAs). In academic settings, digital assistants have evolved from simple keyword-based FAQ bots into context-aware systems using Retrieval-Augmented Generation (RAG).

\subsection{Literature Survey}

\subsubsection{Evolution of Front-End Frameworks}
Modern web development demands fast rendering, modular components, and seamless state management. Frameworks like React have become industry standards due to their virtual DOM reconciliation and component-driven structure. With the release of React 19 and Vite 8, build tools now provide instant hot-module replacement (HMR), reducing bundle sizes and improving loading speed. Tailwind CSS v4 further simplifies styling by offering utility-first classes and CSS variables, allowing developers to create consistent themes without writing heavy external CSS files.

\subsubsection{Multi-Agent Systems in AI}
Traditional AI assistants rely on a single large language model (LLM) prompt to answer all types of queries. However, a single prompt often struggles when switched between creative explanations, strict administrative policy lookups, and technical data extraction. Recent research shows that dividing tasks among a swarm of specialized agents---directed by a central supervisor node---improves retrieval accuracy and reduces incorrect responses. Frameworks like LangGraph allow developers to model agent workflows as directed state graphs, enabling structured decision-making.

\subsection{Technology Stack}
The Chimera AI platform uses a modern tech stack divided cleanly between front-end user experience and back-end processing.

\subsubsection{Front-End Technologies}
\begin{itemize}
    \item \textbf{React 19 \& Vite 8:} Serves as the core front-end framework and lightning-fast bundler for rendering dynamic UI components.
    \item \textbf{Tailwind CSS v4:} Used for layout styling, custom color variables, flexbox/grid alignments, and responsive mobile-first styling.
    \item \textbf{Framer Motion:} Enables clean UI transitions, card hover effects, and smooth page navigation animations.
    \item \textbf{Lucide React Icons \& 21st.dev Components:} Provides UI icons and modern design elements like spotlight cursor cards and glowing ambient backgrounds.
\end{itemize}

\subsubsection{Back-End Technologies}
\begin{itemize}
    \item \textbf{FastAPI (Python):} Asynchronous web framework for handling REST API requests and Server-Sent Events (SSE) token streams.
    \item \textbf{Express API Gateway (Node.js):} Acts as an intermediate gateway running on Port 3000 to manage routing, CORS headers, and client request proxies.
    \item \textbf{LangGraph \& Groq LLM:} Powers the multi-agent decision routing graph and fast language model token generation.
    \item \textbf{Pinecone Vector Database:} Cloud vector store used for indexing and searching PDF syllabus chunks and question bank files.
\end{itemize}
